\documentclass[12pt]{article}
\usepackage[utf8]{inputenc}
\usepackage[T1]{fontenc}
\usepackage[brazil]{babel}
\usepackage[a4paper,margin=2.5cm]{geometry}
\usepackage{enumitem}
\setlist[itemize]{noitemsep,topsep=4pt}
\title{Sprint 3 - S3\_06\\Qualidade de Software e Código Limpo}
\author{Matheus Moreli Durão Silva\\Matrícula 12621CBS217\\Cibersegurança - Bacharelado - Noturno}
\date{2026}
\begin{document}
\maketitle

\section*{Entregável}
Relatório de qualidade, linting e refatoração. Este documento registra o objetivo do card, a solução desenvolvida e a forma de conferência. Os arquivos de código e as evidências citadas permanecem na mesma pasta deste documento.

\section*{O que foi estudado e aplicado}
\begin{itemize}
\item Legibilidade e nomes claros
\item PEP 8
\item Execução de flake8, pylint e black
\item Registro das melhorias antes e depois
\end{itemize}

\section*{Como conferir}
O avaliador pode abrir os arquivos desta pasta, executar os exemplos quando houver código e comparar a saída com a explicação apresentada. Nos cards ligados ao Inventário de Ativos, a solução foi relacionada ao Trabalho 2 sem modificar requisitos já atendidos do CRUD.

\section*{Resultados reais das ferramentas}
Na branch de melhoria, \texttt{python -m flake8} retornou zero ocorrências; \texttt{python -m black .} formatou nove arquivos e \texttt{python -m black --check .} confirmou dez conformes. \texttt{pylint --errors-only} não apontou erros nos módulos centrais; a análise completa atribuiu 9,29/10 e mostrou avisos de muitos argumentos, docstrings e estilo. No exemplo \texttt{depois\_refatoracao.py}, pylint deu 10,00/10.
\section*{Antes, depois e limites}
O exemplo \texttt{antes\_refatoracao.py} tem nomes vagos e formatação fraca; \texttt{depois\_refatoracao.py} usa nomes, validação e docstring. No programa real, operações passaram a métodos de \texttt{Inventario}, persistência a \texttt{RepositorioJSON} e risco a \texttt{AnalisadorRisco}; adaptadores preservam as chamadas antigas. O alerta de muitos atributos decorre dos campos exigidos; as duas lambdas sinalizadas são executadas imediatamente. O relatório \texttt{relatorio\_qualidade.md} registra comandos e análise.
\section*{Conclusão}
Os 45 testes passaram após a formatação. Ferramentas estáticas ajudam a localizar manutenção difícil, mas exigem leitura do código para decidir se cada aviso é defeito.
\end{document}
